\section{Method}
\label{sec:method}

\subsection{Conditional DDPM architecture}
We implement a 1D temporal denoising diffusion probabilistic model with the
following components:

\begin{enumerate}
    \item \textbf{Temporal embedding}: sinusoidal frequencies $f_j = 10000^{-2j/D}$
    embedded into a width-\texttt{width} MLP, where $D$ is half the width.

    \item \textbf{History conditioning}: a separate MLP consumes all 72 past
    residuals, producing a \texttt{width}-dimensional conditioning vector that is
    \emph{added} to the temporal embedding before entering the denoiser.

    \item \textbf{Dilated residual blocks}: four blocks with dilation $\{1, 2, 4, 8\}$,
    each containing GroupNorm (groups=4), SiLU activation, 3-conv padding to
    preserve length, and a linear condition injection.

    \item \textbf{Output head}: GroupNorm + SiLU + 3-conv producing a single-channel
    epsilon prediction at each of the 24 forecast horizons.
\end{enumerate}

Total parameter count: 27,169.

\subsection{Forward and reverse processes}
We implement cosine cumulative noise schedule~\citep{ho2020denoising} with 64
diffusion steps:
\begin{equation}
    \beta_t = 1 - \frac{\cos((t+0.008)/1.008 \cdot \pi/2)^2}{\cos((t-1+0.008)/1.008 \cdot \pi/2)^2},
    \qquad \alpha_t = 1 - \beta_t, \qquad \bar{\alpha}_t = \prod_{i=1}^{t} \alpha_i.
\end{equation}

\subsection{Training}
Training targets \emph{unit Gaussian epsilon} at each timestep. The loss is mean
squared error between the network's epsilon prediction and the actual injected
Gaussian noise:
\begin{equation}
    \mathcal{L} = \mathbb{E}_{t, x_0, \epsilon \sim \mathcal{N}(0,I)} \left[
        \left\| \epsilon - \epsilon_\theta(\sqrt{\bar{\alpha}_t} x_0 + \sqrt{1-\bar{\alpha}_t} \epsilon, t, \text{history}) \right\|^2
    \right].
\end{equation}

Training hyperparameters for the pilot run (Issue \#6): AdamW with $\text{lr}=0.001$,
batch size 128, three epochs, BF16 precision on RTX3090 GPU~4, seed 20261025.

\subsection{Engineering validation guards}
After the pilot revealed catastrophic scale explosion, we added validation guards
that \emph{reject} non-finite or implausibly scaled samples. These are engineering
checks, not model repairs:

\begin{enumerate}
    \item \textbf{Nonfinite rejection}: raise on nonfinite conditioning history,
    nonfinite epsilon prediction at any reverse step, or nonfinite reverse sample
    intermediate state.

    \item \textbf{Train-reference scale screen}: after de-normalization, compare
    sample standard deviation to frozen training target standard deviation. Reject
    if ratio $< 0.1$ (collapse) or $> 5.0$ (dispersion), or if mean offset exceeds
    5 training standard deviations. Both bounds are explicit policy defaults, not
    estimates of physical plausibility.
\end{enumerate}

The historical GPU pilot predates these guards; its samples have SD ratio 22.95
and would be rejected by the screen.
